%%%%%%%%%%%%%%%%%%%%%%%%%%%%%%%%%%%%%%%%%%%%%%%%%%%%%%%%%%%%%%%%%%%%%%
% Disciplined Vibe Coding - A 4-Block Lecture Series
% Author: Prof. Dr.-Ing. Mark Schutera
% Based on Tufte-style layout
%%%%%%%%%%%%%%%%%%%%%%%%%%%%%%%%%%%%%%%%%%%%%%%%%%%%%%%%%%%%%%%%%%%%%%

\documentclass{../sharedAssets/tufte}

%%
% Book metadata
\title{Disciplined Vibe}
\author[Prof. Dr.-Ing. Mark Schutera]{Prof. Dr.-Ing. Mark Schutera}
\publisher[\defaultpublisher]{\defaultpublisher}

%%
% Load optional fonts
\loadoptionalfonts

% Additional packages
\usepackage{pifont}
\usepackage{amsmath}
\usepackage{soul}
\usepackage{textcase}

\begin{document}

% Standard front matter
\makestandardfrontmatter

% Dedication
\makededication{%
  \openepigraph{There's a new kind of coding I call ``vibe coding'', where you fully give in to the vibes, embrace exponentials, and forget that the code even exists.}{Andrej Karpathy, 2025}
  \vspace{2em}
  \openepigraph{Wir wollen offensiv spielen, aber mit h\"ochster Disziplin.}{Joachim L\"ow, 2010}%
}

\tableofcontents

% Introduction
\makeintroduction{%
  These notes accompany a four-block seminar on Disciplined Vibe Coding:
  programming in which the developer steers a Large Language Model in natural
  language while the model writes, edits, and reasons about the code. We take
  Karpathy's framing\sidenote{Coined by Andrej Karpathy in February 2025 and
  building on his earlier essays Software 2.0 and Software
  3.0.\cite{karpathy2025vibe}} as the starting point but ground it in concrete
  craft: the \textsc{skills} catalogue maintained by Matt
  Pocock\cite{pocock2025skills} supplies a working vocabulary of reusable agent
  behaviours, from \texttt{to-prd} to \texttt{tdd} to \texttt{write-a-skill}.\\

  \newthought{The structure} is deliberately interleaved. Each of the four blocks
  opens with a short theoretical introduction and closes with a hands-on
  breakout in which students operate a local model (\textsc{ollama}) through
  the \texttt{opencode} terminal client inside an editor of their choice.\\

  \newthought{These notes are a living thing,} subject to continuous updates.
  They favour breadth and orientation over completeness; the breakouts are
  where the actual learning happens.
}

%% Start the main matter
\mainmatter

%%%%%%%%%%%%%%%%%%%%%%%%%%%%%%%%%%%%%%%%%%%%%%%%%%%%%%%%%%%%%%%%%%%%%%
\chapter{Foundations and the Local Workbench}
\label{ch:foundations}

\section{What Vibe Coding Actually Is}
\label{sec:what-is-vibe}

\newthought{Karpathy's coinage} describes a shift in the locus of authorship.
In Software 1.0, humans wrote explicit instructions; in Software
2.0, gradient descent wrote the weights of neural networks; in Software
3.0\sidenote{Karpathy's Software is Changing (Again) talk frames
\textsc{llm}s as a new computing substrate, ``\textsc{llm}s are the new
operating system'', with prompts as programs and context windows as
\textsc{ram}.\cite{karpathy2025software3}} the program is a natural-language
prompt, and the \textsc{llm} compiles it to behaviour.\index{Software 2.0}\index{Software 3.0}\index{Karpathy}

\begin{quote}
``The hottest new programming language is English.'', Andrej Karpathy
\end{quote}

\newthought{Vibe coding} is the working style that follows from this shift.
The developer states intent, accepts diffs, runs the result, and iterates
without necessarily reading every line. Simon Willison has argued that the
term should be reserved for cases where the code is genuinely not read
by the human, and that most professional \textsc{ai}-assisted development
is something more disciplined.\cite{willison2025notvibe}\index{vibe coding!definition}

This seminar treats vibe coding as a spectrum, parameterised by
how much the human inspects, structures, and constrains the model's
output. \textsc{skills} are the means by which we move along that
spectrum deliberately rather than by accident, hence the qualifier
disciplined.

\newthought{A useful counterweight} to the exponentials rhetoric is
Fred Brooks's No Silver Bullet\cite{brooks1987silver}: the
essential complexity of software, specification, design,
verification, is irreducible, and no tool eliminates it. \textsc{llm}s
attack accidental complexity (typing, boilerplate, look-up); the
essential part still belongs to you. The skills catalogue is, in
effect, a set of disciplines for keeping the essential work in human
hands while delegating the accidental work to the model.\index{Brooks!No Silver Bullet}\index{essential complexity}

\section{Skills as Reusable Prompt Modules}
\label{sec:skills-concept}

\newthought{A skill}, in the Pocock sense, is a small, named, reusable
bundle of instructions and examples that an agent loads on demand, a
``cheat sheet'' in Yegge's terms\cite{yegge2024cheating}, a ``recipe''
in Crawshaw's.\cite{crawshaw2024sketch}\index{skills}\index{prompt engineering}
The underlying principle is the oldest one in the Pragmatic Programmer's
catalogue\cite{hunt1999pragmatic}: \textsc{dry}, Don't Repeat
Yourself. A prompt you have written twice is a skill waiting to be
extracted.\index{DRY}\index{Pragmatic Programmer}

\newthought{Skills are not the only} composable layer. The
Model Context Protocol\sidenote{\textsc{mcp} is an open standard
introduced by Anthropic in late 2024 for connecting \textsc{llm} agents
to external data sources and tools through a uniform
interface.\cite{anthropic2024mcp}} (\textsc{mcp}) provides the
complementary half: where a skill is a prompt module that shapes
how the agent thinks, an \textsc{mcp} server is a tool surface
that determines what the agent can reach, a database, a ticket
system, a filesystem, a browser. Skills and \textsc{mcp}s are
orthogonal; a mature workbench uses both.\index{MCP}\index{Model Context Protocol}

\begin{table*}[ht]
\begin{tabular}{p{0.18\textwidth}p{0.32\textwidth}p{0.34\textwidth}}
\toprule
\textbf{Layer} & \textbf{Artifact} & \textbf{Example} \\
\midrule
Model        & Weights              & Llama 3.2, Qwen 2.5 Coder \\
Harness      & Agent loop, permissions & opencode, Claude Code \\
\textsc{mcp} server & External tool surface & Filesystem, GitHub, Postgres \\
Skill        & Prompt module        & \texttt{tdd}, \texttt{to-prd} \\
Project      & Repository, CLAUDE.md & Your codebase \\
Intent       & The user             & You \\
\bottomrule
\end{tabular}
\vspace{1em}
\caption{The six layers of a vibe-coding stack.}
\label{tab:layers}
\end{table*}

\section{Breakout 1: Setting Up the Local Workbench}
\label{sec:breakout-1}

\newthought{The goal} of this breakout is a fully local, offline-capable
vibe-coding environment: an \textsc{ide}, a local model served by
\textsc{ollama}, and the \texttt{opencode} terminal client wired to it.

\subsection*{Step 1: Install an Editor}

Vibe coding can in principle be done from a bare terminal, but an editor
that renders diffs and shows the file tree makes inspection radically
easier.\marginnote{The default in our bootstrap script is \texttt{micro},
a $\sim$10\,\textsc{mb} terminal editor with mouse support and no learning
curve, \texttt{Ctrl-S} to save, \texttt{Ctrl-Q} to quit. It keeps the whole
workbench (editor, opencode, shell) in one window. Heavier alternatives,
Zed, Sublime, \textsc{vs code}, JetBrains, all work equally well; the point
is: you should be able to see what the agent changes.}

\begin{fullwidth}
\begin{verbatim}
# macOS
brew install micro

# Linux (Debian/Ubuntu)
sudo apt-get install micro
\end{verbatim}
\end{fullwidth}

\marginnote{The companion \texttt{setup/bootstrap.sh} script in the lecture
repository performs all four steps of this breakout in one shot, with
\texttt{EDITOR=micro|zed|vscode|none} and \texttt{MODEL=...} as override
knobs. Re-runs are idempotent.}

\subsection*{Step 2: Install Ollama and Pull a Coding Model}

\textsc{ollama}\cite{ollama2024} is a local model runner that exposes an
\textsc{http} \textsc{api} compatible with much of the OpenAI surface.
On a machine with \(\geq 16\,\text{GB}\) of \textsc{ram} a 7B coder model
runs comfortably; on \(\geq 32\,\text{GB}\) you can host a 14B model.

\begin{fullwidth}
\begin{verbatim}
# Install Ollama
brew install ollama          # macOS
curl -fsSL https://ollama.com/install.sh | sh   # Linux

# Start the server (runs on http://localhost:11434)
ollama serve

# Pull a coding-tuned model in another terminal
ollama pull qwen2.5-coder:7b
ollama run  qwen2.5-coder:7b "write a python fizzbuzz"
\end{verbatim}
\end{fullwidth}

\marginnote{If \texttt{qwen2.5-coder} is too heavy, fall back to
\texttt{llama3.2:3b} or \texttt{deepseek-coder:6.7b}. The skill of
choosing the smallest model that does the job is itself part of
vibe coding.}

\subsection*{Step 3: Install opencode and Connect It to Ollama}

\texttt{opencode}\cite{opencode2025} is a terminal-native agent harness
modelled after Claude Code. It can speak to any OpenAI-compatible
endpoint, which is exactly what \textsc{ollama} provides.

\begin{fullwidth}
\begin{verbatim}
# Install opencode
curl -fsSL https://opencode.ai/install | bash

# Configure a local provider (~/.config/opencode/config.json)
{
  "provider": {
    "ollama": {
      "npm": "@ai-sdk/openai-compatible",
      "options": { "baseURL": "http://localhost:11434/v1" },
      "models": { "qwen2.5-coder:7b": { "name": "Qwen 2.5 Coder 7B" } }
    }
  }
}

# Launch opencode in your project directory
cd ~/projects/playground && opencode
\end{verbatim}
\end{fullwidth}

\subsection*{Step 4: Smoke Test}

Inside \texttt{opencode}, ask: ``create a Python script that prints
the first ten Fibonacci numbers, then run it.'' If you see a diff in your
\textsc{ide} and the output \texttt{0 1 1 2 3 5 8 13 21 34} in the
terminal, the workbench is live.

\vspace{1em}
\dots

%%%%%%%%%%%%%%%%%%%%%%%%%%%%%%%%%%%%%%%%%%%%%%%%%%%%%%%%%%%%%%%%%%%%%%
\chapter{Planning and Design Skills}
\label{ch:planning}

\section{Why Plan Before You Vibe}
\label{sec:why-plan}

\newthought{The single largest failure mode} of naive vibe coding is
premature implementation: the model happily produces 400 lines of
code for a problem the user has not yet specified. The planning skills in
the Pocock catalogue exist to slow that reflex down.\index{planning}

Addy Osmani has named this pattern the ``70-percent problem'':
\textsc{ai}-assisted code feels 70\% done almost immediately, and the
remaining 30\%, the part that requires actual
specification, becomes disproportionately
painful.\cite{osmani2024seventy}\index{70-percent problem}
Eric Evans's older prescription points to the same root cause: agree on
the ubiquitous language before you write
code.\cite{evans2003ddd}\index{ubiquitous language}\index{DDD}

\section{The Five Planning Skills}
\label{sec:planning-skills}

\begin{margintable}
\caption{Planning skills in the Pocock catalogue.}
\label{tab:planning-skills}
\vspace{1em}
\begin{tabular}{ll}
\toprule
\textbf{Skill} & \textbf{Output} \\
\midrule
\texttt{to-prd}              & GitHub issue (PRD) \\
\texttt{grill-me}            & Resolved decision tree \\
\texttt{design-an-interface} & N parallel designs \\
\texttt{request-refactor-plan} & Tiny-commit plan \\
\texttt{to-issues}           & Vertical-slice tickets \\
\bottomrule
\end{tabular}
\end{margintable}

\newthought{\texttt{to-prd}} converts an open-ended chat into a Product
Requirements Document filed as a GitHub issue. The point is not the
document but the forcing function: writing a PRD reveals what you
do not yet know.\index{to-prd}

\newthought{\texttt{grill-me}} interviews the user relentlessly about a
plan, walking each branch of the decision tree until every fork is
either closed or explicitly deferred. This is the conversational analog
of Kent Beck's ``make the change easy, then make the easy
change''\cite{beck2024tidy}.\index{grill-me}

\newthought{\texttt{design-an-interface}} spawns parallel sub-agents that
each propose a different shape for the same module, a deliberate
exercise in designing it twice. The cost of generating three
interfaces with an \textsc{llm} is minutes; the cost of choosing the
wrong one in production is months.\index{design-an-interface}

\newthought{\texttt{request-refactor-plan}} produces a sequence of atomic
commits, each independently reviewable. This maps directly onto Beck's
tidying-first discipline, onto Fowler's refactoring catalogue, small,
behaviour-preserving transformations under green
tests\cite{fowler2018refactoring}, and onto the trunk-based development
literature.\index{refactor}\index{Fowler}

\newthought{\texttt{to-issues}} encodes Hunt and Thomas's
tracer-bullet doctrine\cite{hunt1999pragmatic}: rather than
specifying the whole system up front, ship a thin end-to-end slice that
exercises every layer, then iterate. Vertical slicing is what makes a
PRD survive contact with reality, and it is the failure mode most
commonly skipped by an over-eager agent.\index{tracer bullets}\index{vertical slicing}

\section{Breakout 2: Plan a Feature Without Writing It}
\label{sec:breakout-2}

\newthought{The exercise.} Pick a small feature for a personal project, a
\textsc{cli} that converts \textsc{json} to Markdown tables, say, and
walk it through four skills back-to-back, in this order:

\begin{enumerate}
  \item \texttt{grill-me}: let the agent interview you until it has no
        more questions.
  \item \texttt{to-prd}: emit the resulting PRD as a GitHub issue
        (or a local Markdown file if you are offline).
  \item \texttt{design-an-interface}: generate three competing
        \textsc{api} shapes; pick one and write a one-paragraph
        rationale.
  \item \texttt{to-issues}: decompose the chosen design into
        vertical slices.
\end{enumerate}

\marginnote{Constraint: do not let the model write any implementation
code during this breakout. The deliverable is a plan you trust,
not a running program.}

\vspace{1em}
\dots

%%%%%%%%%%%%%%%%%%%%%%%%%%%%%%%%%%%%%%%%%%%%%%%%%%%%%%%%%%%%%%%%%%%%%%
\chapter{Development Skills}
\label{ch:development}

\section{Tests as the Vibe-Coder's Compass}
\label{sec:tdd-theory}

\newthought{Test-Driven Development}, due to Kent Beck\cite{beck2002tdd},
predates \textsc{llm}s by two decades but maps onto vibe coding with
unusual precision. The red-green-refactor loop gives the agent an
unambiguous target, a failing test, and an unambiguous stopping
condition, a passing one.\index{TDD}\index{red-green-refactor}

\[
\underbrace{\text{red}}_{\text{spec}} \;\longrightarrow\;
\underbrace{\text{green}}_{\text{minimum impl.}} \;\longrightarrow\;
\underbrace{\text{refactor}}_{\text{shape}}
\]

\marginnote{This loop also serves as a hallucination filter: the
test either passes or it does not. The agent cannot vibe its way past a
red bar.}

\section{Root-Cause over Symptom}
\label{sec:triage-theory}

\newthought{\texttt{triage-issue}} encodes the discipline of finding the
true cause of a bug before patching its symptom. This is the same
medicine that Toyota's five whys prescribes and that
\textsc{sre} post-mortems formalise. Vibe coding without it produces
what Crawshaw calls ``mole-whacking''\cite{crawshaw2024sketch}: each
patch creates two new bugs.\index{triage}\index{root cause}

\newthought{Michael Feathers} gives the operational recipe in
Working Effectively with Legacy Code\cite{feathers2004legacy}:
before changing behaviour, write a characterization test that
pins down current behaviour, then locate a seam at which the
change can be made safely. \texttt{triage-issue} pushes the agent down
exactly this path, first reproduce, then isolate, then fix.\index{Feathers}\index{characterization tests}\index{seams}

\section{Architecture Through Domain Language}
\label{sec:architecture-theory}

\newthought{\texttt{improve-codebase-architecture}} reads two artefacts
specifically: a \texttt{CONTEXT.md} that names the domain concepts, and
the \texttt{docs/adr/} folder that records past decisions. Together
these are the project's ubiquitous language and its
architectural memory, the two inputs Evans argues are
prerequisite to any non-trivial refactor.\cite{evans2003ddd}

\newthought{John Ousterhout's} A Philosophy of Software
Design\cite{ousterhout2018philosophy} sharpens the criterion: prefer
deep modules, narrow interfaces hiding substantial
implementation, over wide, shallow ones. When \texttt{design-an-interface}
emits three competing shapes, the question to ask is not ``which is
prettiest'' but ``which hides the most complexity behind the smallest
surface.''\index{Ousterhout}\index{deep modules}

\begin{table*}[ht]
\begin{tabular}{p{0.18\textwidth}p{0.36\textwidth}p{0.32\textwidth}}
\toprule
\textbf{Skill} & \textbf{When to reach for it} & \textbf{Anchor concept} \\
\midrule
\texttt{tdd}            & New behaviour, fixable bug & Beck, red-green-refactor \\
\texttt{triage-issue}   & Reported defect            & Five whys, root cause \\
\texttt{improve-...-architecture}  & Drift, friction & Evans, ubiquitous language \\
\texttt{scaffold-exercises} & Teaching, demos & Worked example pedagogy \\
\bottomrule
\end{tabular}
\vspace{1em}
\caption{Development skills mapped to their classical anchors.}
\label{tab:dev-skills}
\end{table*}

\section{Reading Diffs and Resolving Conflicts}
\label{sec:diffs-merges}

\newthought{Inspecting what the agent changed} is the line between vibe
coding and engineering. The interface for that inspection is the
\emph{unified diff}, the same format that \texttt{git diff} prints and that
\texttt{opencode} renders before applying any edit. Reading diffs is a
prerequisite skill, not an advanced one.\index{diff}\index{unified diff}

\begin{fullwidth}
\begin{verbatim}
--- a/greet.py
+++ b/greet.py
@@ -1,2 +1,2 @@
-def greet(name):
-    return "Hello, " + name + "!"
+def greet(name, lang="en"):
+    return f"Hello, {name}!"
\end{verbatim}
\end{fullwidth}

\marginnote{The hunk header \texttt{@@ -1,2 +1,2 @@} reads as: ``the old
file's hunk starts at line 1 and is 2 lines long; the new file's hunk
starts at line 1 and is 2 lines long.'' Diffs are read in pairs of
\texttt{-}/\texttt{+} blocks, not top-to-bottom.}

\newthought{A merge conflict} happens when two streams of changes touch
the same line and \textsc{git} cannot decide which to keep. It writes
both versions into the file with markers and hands the decision back
to the human:\index{merge conflict}

\begin{fullwidth}
\begin{verbatim}
<<<<<<< HEAD
return f"Hello, {name}!"
=======
greeting = {"en": "Hello", "de": "Hallo"}[lang]
return greeting + ", " + name + "!"
>>>>>>> feature/i18n
\end{verbatim}
\end{fullwidth}

\noindent
The markers delimit \emph{ours} (\texttt{HEAD}, the branch you are on)
and \emph{theirs} (the branch being merged in). Three resolution
strategies, in increasing order of judgement required:

\begin{enumerate}
  \item Take ours wholesale: \texttt{git checkout --ours FILE}.
  \item Take theirs wholesale: \texttt{git checkout --theirs FILE}.
  \item Hand-merge: open the file, delete the markers, and write the
        union that captures both intents.
\end{enumerate}

\marginnote{Most real conflicts demand the third option. The first two
are tempting precisely because they are fast, and wrong precisely because
they silently discard work, the merge-time analogue of
\texttt{rm -rf}.}

\newthought{The three-way merge} is the conceptual model behind
\texttt{git diff --base}. A clean resolution is a function of three
inputs: \emph{ours}, \emph{theirs}, and the common ancestor (\emph{base}).
Looking at \emph{ours} versus \emph{theirs} alone is solving the problem
with one eye closed; the base tells you which side actually changed
what.\index{three-way merge}

\begin{fullwidth}
\begin{verbatim}
git diff --ours   FILE   # what you had before the merge
git diff --theirs FILE   # what is being merged in
git diff --base   FILE   # vs. the common ancestor
git merge --abort        # escape hatch: throws the merge away
\end{verbatim}
\end{fullwidth}

\newthought{Three vibe-coding-specific failure modes} around diffs:
(i)~\emph{accepting a hunk without reading the header}, the agent removed
a function call three lines above the visible \texttt{+} and you missed
it; (ii)~\emph{letting the model resolve the conflict for itself},
conflicts encode intent collisions that the model has no way to
adjudicate; (iii)~\emph{committing with conflict markers still in the
file}, which a pre-commit hook would have caught (cf.\ \S\ref{sec:guardrails}).

\marginnote{The companion script \texttt{setup/diffs\_and\_merges.sh}
builds a sandbox repository at \texttt{/tmp/vibe-merge-demo} with a
deliberately engineered conflict between an \texttt{f}-string refactor
and an \texttt{i18n} parameter addition. Both changes are correct in
isolation; only the hand-merge is correct overall. The accompanying
\texttt{diffs\_and\_merges.md} explains every step in prose.}

\section{Breakout 3: Fix a Real Bug, TDD-Style}
\label{sec:breakout-3}

\newthought{The exercise.} Take an existing open-source repository
(or one of your own) with a known issue. With \texttt{opencode}
pointed at your local model:

\begin{enumerate}
  \item Invoke \texttt{triage-issue}: have the agent locate the
        root cause and propose a TDD plan.
  \item Invoke \texttt{tdd}: write the failing test first,
        confirm red, then let the agent close it.
  \item Inspect the diff before accepting. This is the moment
        where vibe coding becomes engineering.
\end{enumerate}

\marginnote{Karpathy's framing reminds us to ``embrace exponentials''
\cite{karpathy2025vibe}; Beck's framing reminds us that an exponential
without a test is just a faster way to be wrong.}

\vspace{1em}
\dots

%%%%%%%%%%%%%%%%%%%%%%%%%%%%%%%%%%%%%%%%%%%%%%%%%%%%%%%%%%%%%%%%%%%%%%
\chapter{Tooling, Guardrails, and Authoring Your Own Skill}
\label{ch:tooling}

\section{The Harness Is Part of the Program}
\label{sec:harness}

\newthought{A vibe-coding session} is not just a model and a prompt; it
is a harness with hooks, permissions, and a working directory. Every
file written, every command run, and every git operation passes through
that harness. Configuring it deliberately is the difference between a
toy and a tool.\index{harness}\index{hooks}

\newthought{Tool access} flows through \textsc{mcp} servers attached to
the harness. A minimal \texttt{opencode} configuration that wires in a
filesystem server and a GitHub server looks as follows:

\begin{fullwidth}
\begin{verbatim}
# ~/.config/opencode/config.json
{
  "mcp": {
    "filesystem": {
      "command": "npx",
      "args": ["-y", "@modelcontextprotocol/server-filesystem", "."]
    },
    "github": {
      "command": "npx",
      "args": ["-y", "@modelcontextprotocol/server-github"],
      "env":  { "GITHUB_TOKEN": "${env:GITHUB_TOKEN}" }
    }
  }
}
\end{verbatim}
\end{fullwidth}

\marginnote{Add only the \textsc{mcp} servers you actually need. Each
server expands the agent's blast radius; an unused \texttt{shell-exec}
\textsc{mcp} is a liability, not a capability.}

\section{Pre-Commit Hooks and Git Guardrails}
\label{sec:guardrails}

\newthought{\texttt{setup-pre-commit}} installs Husky with lint-staged,
type-checking, and tests. The agent now cannot land a commit that
breaks the build, a hard constraint that prevents most ``70-percent''
regressions before they enter \textsc{git} history.\index{pre-commit}

\newthought{\texttt{git-guardrails-claude-code}} goes further and blocks
genuinely destructive commands (\texttt{push --force},
\texttt{reset --hard}, \texttt{branch -D}) at the harness layer.
Willison's writing on prompt injection\cite{willison2024injection}
is the conceptual backdrop: an agent that can be steered by untrusted
input must not also have unrestricted access to your remote.\index{prompt injection}\index{guardrails}

\section{Writing Your Own Skill}
\label{sec:write-a-skill}

\newthought{\texttt{write-a-skill}} is the meta-skill. It enforces
progressive disclosure: the skill's top-level description must be
short enough to fit in the agent's selection prompt, while the body and
bundled resources are loaded only when the skill is actually invoked.
This is the same design principle as Unix \texttt{man} pages, and
the reason a catalogue of dozens of skills can coexist without flooding
the context window.\index{progressive disclosure}\index{write-a-skill}

\newthought{The granularity rule} for writing a good skill mirrors
Robert C.\ Martin's rule for writing a good function in
Clean Code\cite{martin2008clean}: it should do one thing,
do it well, and be nameable in a verb-phrase. \texttt{tdd} passes;
\texttt{do-stuff-with-tests} does not. McConnell's
Code Complete\cite{mcconnell2004complete} adds the missing
ingredient: a checklist. Pre-commit hooks and git guardrails are
checklists rendered executable, the harness will not let the agent
forget what you would otherwise forget under pressure.\index{Clean Code}\index{Code Complete}\index{checklists}

\begin{fullwidth}
\begin{verbatim}
, 
name: my-domain-glossary
description: Extract a DDD-style ubiquitous-language glossary from
             the current conversation and write it to CONTEXT.md.
, 

# Body
1. Scan the current transcript for domain nouns and verbs.
2. Group them by aggregate.
3. Emit a Markdown table to CONTEXT.md.
\end{verbatim}
\end{fullwidth}

\section{Breakout 4: Author and Use Your Own Skill}
\label{sec:breakout-4}

\newthought{The exercise.} Identify one task you have done by hand at
least three times this term, tagging exam questions, summarising
seminar notes, drafting issue triage comments, and convert it into a
skill.

\begin{enumerate}
  \item Run \texttt{write-a-skill} and let the agent scaffold the
        directory layout.
  \item Edit the description until it would unambiguously match
        your task and not someone else's.
  \item Install it (\texttt{\textasciitilde/.config/opencode/skills/}
        or the equivalent for your harness).
  \item Invoke it on a fresh instance of the task. Iterate until the
        output is acceptable without manual edits.
  \item Configure \texttt{setup-pre-commit} and
        \texttt{git-guardrails-claude-code} on the same project so the
        skill cannot regress your repository.
\end{enumerate}

\marginnote{The deliverable is a \texttt{.zip} of your skill folder
plus a one-page reflection: what was hard to specify, and why?}

\vspace{1em}
\dots

%%%%%%%%%%%%%%%%%%%%%%%%%%%%%%%%%%%%%%%%%%%%%%%%%%%%%%%%%%%%%%%%%%%%%%
\chapter*{Closing Notes}
\addcontentsline{toc}{chapter}{Closing Notes}

\newthought{Vibe coding} is neither magic nor a substitute for
engineering judgement. It is a new interface to the same craft,
and like any interface it rewards practitioners who understand the
layers underneath. The four blocks of this seminar trace those layers
from the model on disk to the skill in your editor; the breakouts are
where the understanding actually lands.

\begin{quote}
``\textsc{ai}-assisted programming will not replace the deliberate
practice of software engineering. It will, however, make the
practitioners who refuse to develop that practice obsolete much
faster.'', paraphrasing Crawshaw\cite{crawshaw2024sketch}
\end{quote}

\makestandardbackmatter{vibecoding_references}

\end{document}
